\section{Storage Architecture}

\label{sec:storage}

\subsection{Three-Layer Design}

The Experience Ledger uses a three-layer storage architecture:

\begin{table}[t]
\centering
\begin{tabular}{lccc}
\toprule
\textbf{Layer} & \textbf{Purpose} & \textbf{Latency} & \textbf{Cost} \\
\midrule
On-chain & Root verification & 1-10s & High \\
\IPFS{} & Active storage & 100-500ms & Low \\
Arweave & Permanent archive & 1-5s & Medium \\
\bottomrule
\end{tabular}
\caption{Storage layer characteristics}
\label{tab:storage}
\end{table}

\subsection{\IPFS{} Integration}

Experiences are stored as IPLD (InterPlanetary Linked Data) nodes:

\begin{lstlisting}[caption=IPLD Experience Node]
{
  "Data": <CBOR-encoded experience>,
  "Links": [
    {"Name": "embedding", "Cid": "bafy..."},
    {"Name": "metadata", "Cid": "bafy..."}
  ]
}
\end{lstlisting}

\subsection{Arweave Permanence}

Critical library snapshots are archived to Arweave:

\begin{equation}
\text{Cost}_{\text{Arweave}}(n) = 0.001 \times n \text{ bytes} \times \text{AR/GB}
\end{equation}

At current rates (\$15/GB for 200+ years), archiving 10,000 experiences costs approximately \$0.50.

\subsection{Content Deduplication}

Content-addressing provides automatic deduplication:

\begin{theorem}[Deduplication Savings]
For $n$ experiences with $d$ duplicates, storage cost is $O((n-d) \cdot s)$ where $s$ is average experience size.
\end{theorem}

In practice, we observe 15-30\% deduplication across experience libraries due to common reasoning patterns.

